<<P000907> Импорт
<<P000885>

<<P000773>%
<<<P000614> Библиотека, чье экспортируемое пространство имен
(или подмножество его отображений) доступно в текущей библиотеке.

<<P000485>
<libraryImport> ::= <metadata> <importSpecification>

<importspecification> ::= <<<P000908><
<<<P000909>< <configurableUri>
(<<P000910>). \AS{} <typeIdentifier>)?
<комбинатор>*; "
<<P000509>

<<P000774>%
Интерпретация настраиваемых URI описана в другом месте.
(<<P000533>).
Импорт определяет URI <<<P000005>
где находится декларация импортируемой библиотеки.
Ошибка времени компиляции, если указанный URI импорта
Это не относится к библиотечной декларации.

<<<P000775>%
<<<P000615> В настоящее время ведется работа над библиотекой.
Импорт изменяет пространство имен текущей библиотеки.
в порядке, определяемом импортной библиотекой и
факультативные элементы импорта.

<<P000776>%
Импорт может быть отложен или немедленным.
А.
\IndexCustom{deferred import}{import!deferred}
Отличается возникновением
встроенный идентификатор <<<P000913>{} после URI.
Ан
<<<P000914>{немедленный импорт}{import!immediate}
Это импорт, который не откладывается.

<<P000777>%
Немедленная директива по импорту <<<P000006> может необязательно включать
а <<<P000616> формы <<<P000646>>, используемой для префикса
Имя привезено по адресу <<<P000007>.
В данном случае мы говорим, что <<P000753> <<<P000617>,
Или просто «P000618».

<<<P000915>{%>
Обратите внимание, что грамматика обеспечивает отсроченный импорт.
включает в себя пункт префикса,
Таким образом, мы можем ссылаться на пункт префикса <<<P000916>{the} отложенного импорта.%
?

<<<P000778>%
Это ошибка времени компиляции, если префикс используется в отложенном импорте.
Он также используется в качестве приставки к другому пункту импорта.
Ошибка времени компиляции <<<P000754> является префиксом импорта,
и текущая библиотека объявляет члена верхнего уровня с именем базы <<<P000917>.

<<<P000779>%
Директива по импорту <<<P000008> может необязательно включать положения комбинатора пространства имен
Используется для ограничения набора имен, импортируемых по <<<P000009>.
Их синтаксис, использование и влияние на пространства имен описаны в другом месте.
(<<P000534>, <<P000535>,
<<P000536>.

<<<P000780>%
Библиотека ядра дротика <<<P000647>
Неявно вносится в каждую библиотеку дарта, кроме самой себя.
через импортную оговорку формы
<<P000918>>>{<<P000919>>>{} 'dart:core;}
Если библиотека не импортирует <<P000648>.
любой импорт <<<P000649>,
Если вы ограничены через <<P000920>, <<<P000921> или <<P000922>,
Предотвращает автоматический импорт.

<<<P000923><%
Было бы неплохо, если бы в P000650 не было ничего особенного.
Однако его использование является распространенным,
Это приводит к решению импортировать его автоматически.
С другой стороны, некоторые библиотеки <<<P000010> может пожелать определить
с именами, используемыми <<<P000651>
(что легко сделать, поскольку имена, объявленные библиотекой, имеют приоритет).
Другие библиотеки могут пожелать использовать <<<P000011>,
и может использовать членов <<<P000012> Это противоречит основной библиотеке.
без использования префикса и без возникновения ошибок.
Вышеприведенное правило делает это возможным.
по существу отменить <<<P000652> Специальный режим
С помощью еще одного специального правила.%
?

<<P000924> Импортируемое пространство имен
<<P000886>

<<<P000781>%
Ниже мы указываем импортное пространство имен библиотеки <<<P000013>,
<<P000925>>>{<<P000926>>>{import}}
и использовать это для определения пространства имен, которое определяет объем библиотеки <<<P000014>.

<<P000782>%
Мы должны ввести системные библиотеки, потому что у них есть особые правила.
A <<<P000619> - это библиотека, которая является частью реализации Dart.
Любая другая библиотека представляет собой <<<P000620>.

<<<P000927>{%>
Системную библиотеку можно распознать, если
URI, начинающийся с <<<P000653>.%
?

<<<P000928>{%>Особые правила для системных библиотек существуют по следующей причине.
Нормальные конфликты решаются во время развертывания.
Функциональность системной библиотеки
вводится в приложение во время выполнения,
Они могут меняться с течением времени по мере обновления платформы.
Возможны конфликты с системной библиотекой.
вне контроля застройщика.
Чтобы избежать взлома развернутых приложений,
Конфликты с системными библиотеками рассматриваются специально.
?

<<P000783>%
<<<P000015> Будьте импортной директивой и позвольте <<<P000016> Библиотека, импортированная по адресу <<<P000017>.
The
<<<P000929>{пространство имен, предоставленное}{%
Пространство имен! в соответствии с импортной директивой
Директива по импорту <<<P000018> Пространство имен, полученное при применении
Комбинаторы пространства имен <<<P000019>
в экспортируемое пространство имен <<<P000020>
(P000537).

<<<P000930><%
Обратите внимание, что пространство имен предусмотрено директивой импорта <<<P000021> не является
То же самое, что и пространство имен, импортируемое из <<<P000022>.
Последнее включает в себя разрешение конфликтов.
и определено ниже в этом разделе.%
?

<<P000784>%
Пусть \NamespaceName{<<P000932>>>{local}} будет локальным пространством имен $L$
(P000538).
Пусть <<P000933>{I}{1}{m} будут директивами импорта <<<P000024>,
И пусть <<<P000934>{L}{1}{m} быть библиотеками
Например, <<<P000025> <<<P000026> для всех <<<P000027>.
<<<P000935>{%>
Не является ошибкой иметь многократный импорт одной и той же библиотеки.
?

<<<P000785>%
Пусть <<<P000028>.

<<<P000786>%
<<P000936> Шаг первый.
На первом этапе мы вычисляем пространство имен, полученное из каждой импортируемой библиотеки.
соответственно каждой группе библиотек, импортируемых с одинаковым префиксом,
<<P000937>>>{<<P000938>>>{one},i}.

<<<P000787>%
<<P000939> Шаг первый для импорта без приставки
Когда <<<P000029> не имеет префикса, \NamespaceName{<<P000941>>>{one},i}
из пространства имен, предусмотренного директивой импорта <<<P000030>
Удаление всех привязок к имени, базовое имя которого
то же самое, что и базовое имя декларации верхнего уровня в <<<P000031>,
или имя основания которого является префиксом импортной директивы в <<<P000032>.

<<<P000942>{%>
Этот шаг гарантирует, что \SHOW{} и \HIDE{} принимаются во внимание директивы,
Префикс импорта, а также местная декларация
Импортное имя.%
?
<<P000945>

<<P000788>%
<<P000946> Шаг первый для префиксированного импорта
На этом этапе мы решаем конфликты имен.
Среди названий, импортируемых с тем же приставкой.
Когда <<P000033> имеет приставку <<<P000034>,
давать \List{I'}{1}{k} является подсписком <<P000948>>>{I}{1}{m} с префиксом <<P000035>,
И пусть <<P000949> L'{1}{k} - соответствующие библиотеки.
<<P000950>>>{(который является подсписком <<P000951>>>{L}{1}{m})}.
Пусть \NamespaceName{<<P000953>>>{экспорт},j}, <<P000036>,
быть пространством имен, предусмотренным директивой импорта <<<P000037>
<<<P000954>{, которое \SHOW{} и <<P000956>>>{} в расчете)}.

<<<P000789>%
Когда <<P000038> является неотложенным импортом:
Пусть <<P000957>>>{<<P000958>>>{prefix},i} будет пространством имен, полученным при применении
Конфликт сливается в
<<P000959><<P000960> NS}}{\metavar{exported},1}{<<P000962>>>{exported},k}
(<<P000539>).

<<<P000790>%
Когда <<<P000039> Отсроченный импорт:
В данном случае <<<P000040> 1.
<<<P000963>{(в противном случае произошла бы ошибка времени компиляции)},
и \NamespaceName{<<P000965>>>{prefix},i} - пространство имен, полученное
Добавление связывания имени <<P000654>
неявно вызванное объявление функции с подписью
<<P000655>
<<P000966>{<<P000967>{экспортировано},1}.

<<<P000791>%
Тогда \NamespaceName{<<P000969>>>{one},i} является пространством имен, которое имеет
одно связывание, которое отображает <<<P000041> в \NamespaceName{\metavar{prefix},i}.

<<<P000792>%
В этой ситуации мы говорим, что \NamespaceName{\metavar{one},i} является
<<P000621>,
Потому что он отображает префикс библиотеки в пространстве имен.
<<<P000974>{%>
Пространство имен префикса не является объектом Дарта.
Это просто устройство, которое используется для управления
Выражения формы <<<P000975>{qualifiedName} и что они означают
во время статического анализа.
Следовательно, любая попытка использовать префикс пространства имен в качестве объекта
Ошибка времени компиляции, например,
Это не может быть результатом оценки выражения.

Обратите внимание, что если <<<P000042> и <<<P000043> являются такими, что
<<<P000044> и <<P000045> Оба имеют одинаковую приставку <<<P000046>
затем
$<<P000976>>>{<<P000977>>>{one},l} = <<P000978>>>{<<P000979>>>{one},q} =
<<P000980>>>{<<P000981>>>{one},l} <<P000982> <<P000983>{<<P000984>>>{one},q}$%
?
<<P000985>

<<<P000793>%
<<P000986> Шаг второй.
На втором этапе мы решаем конфликты на высшем уровне.
среди пространств имен, полученных на первом этапе.

<<<P000794>%
<<<P000622> <<<P000047>, \NamespaceName{<<P000988>>>{import}}
Это результат применения конфликта, сливающегося с пространствами имен.
<<P000989><<P000990> NS}}{\metavar{one},1}{<<< P000992>>>{one},m}.

<<<P000795>%
<<<P000048> быть набором заявлений о продлении
(<<P000540>)
со следующими членами:
Декларация о продлении <<<P000049> является членом <<<P000050>
Если есть <<<P000051> такой, что
<<P000052> находится в пространстве имен, предусмотренном директивой импорта <<<P000053>
<<<P000993>{%>
(обратите внимание, что это <<P000994>>>{} и <<P000995>>>{} с учетом,
и включает в себя расширения, импортированные как без, так и с префиксом
?
<<<P000054> В действующей библиотеке не указан
<<P000996>{(что может иметь место, если он импортирует сам)}.
Пусть <<P000997>>>{<<P000998>>>{расширения}} будет пространством имен, которое
для каждого расширения <<<P000055> <<<P000056> На карте новое имя <<<P000057>.

<<<P000999>{%>
\NamespaceName{\metavar{расширения}} дает новое имя
предоставление имплицитного доступа к каждому расширению, экспортируемому импортируемой библиотекой;
и не удаляется \HIDE{{} или <<<P001003>,
Даже те, которые не могут быть доступны, используя их заявленное имя,
Из-за столкновения имени.%
?

<<P000796>%
<<<P000623> <<<P000058> тогда
$\NamespaceName{\metavar{local}}\cup
\NamespaceName{<<P001008>>>{import}}\cup
\NamespaceName{<<P001011>>>{расширения}}$.

<<P000797>%
Пусть <<<P000059>.
Если <<<P000060> импортируется без приставки,
<<<P001012>{именное пространство импортировано из}{библиотека! Именная область импортируется из
<<P000061> Конфликт сузил пространство имен
(<<P000541>)
\NamespaceName{\metavar{one},i}.
%
Иначе <<P000062> импортируется с приставкой <<<P000063>,
В этом случае <<<P001015>{именное пространство импортировано из} <<P000064>
Конфликт сузил пространство имен
<<<P001016>{<<P001017>{экспортировано},i}.

<<<P001018>{%>
Имя пользователя: <<P000065> содержит привязки, экспортируемые <<<P000066>,
За исключением тех, которые удалены комбинаторами пространства имен,
За исключением тех, которые были устранены путем слияния конфликтов.%
?

<<P000798>%
<<<P000067> Библиотека с импортным пространством имен <<<P001019>.
Мы говорим, что имя
\IndexCustom{imported by $L$}{imported!name)
Если имя является ключом <<<P001021>.
Мы говорим, что декларация является
\IndexCustom{imported by $L$}{imported!declaration}}
если декларация является значением <<<P001023>.
Мы говорим, что имя
<\IndexCustom<<<P000070> с префиксом $p$}{imported!name, с приставкой}
если имя является ключом <<<P001025>{}{p}.
Мы говорим, что декларация является
<<<P001026><<<P000072> с префиксом <<<P000073>{%}
Импорт! Декларация с префиксом
если декларацией является значение <<<P001027>{}{p}.


<<<P001028>{Семантика импорта}
<<P000887>

<<P000799>%
<<<P000074> быть директивой импорта, которая относится к URI через строку <<<P000075>.
Семантика <<<P000076> Указывается следующее:

<<<P000800>%
<<<P001029> Семантика отложенного импортаЕсли <<P000077> является отложенным импортом с префиксом <<<P000078>, связыванием <<<P000079> до
<<P000624>
\NamespaceName{\metavar{deferred}}
Присутствующее в библиотеке пространство имен текущей библиотеки <<P000080>.
Пусть \NamespaceName{\metavar{import},i}
пространство имен, импортируемое из библиотеки, указанной в <<<P000081>,
как определено ранее
(<<P000542>).
\NamespaceName{\metavar{deferred}} имеет следующие связывания:

<<P000486>
<<<P001036> Оригинальное название: P000656 привязанный к
Функция с подписью <<P000657>.
Эта функция возвращает будущее <<<P000082>.
При вызове функция вызывает
Немедленный импорт <<P000083> быть исполнены в будущем,
где <<P000084> Происходит от <<<P000085> Удаление слова <<<P001037>
Добавить <<<P001038> <<P000658> Комбинаторное предложение.
Выполнение немедленного импорта может привести к неудаче.
по конкретным причинам реализации.
<<<P001039>{%>
Например, <<P000086> Библиотека отличается от той, что
что указанный URI упоминается во время компиляции;
или ошибка чтения файла уровня ОС; и т.д.%
?
Мы говорим, что вызов <<<P000659>
\IndexCustom{succeeds}{loadLibrary!succeeds} если $f$ завершается ценностью,
и что призыв
\IndexCustom{fails}{loadLibrary!fails} если $f$ завершается ошибкой.
<<P001042>
Для каждой функции верхнего уровня <<<P000089> Оригинальное название: P000755 в
\NamespaceName{<<P001044>>>{import},i}
Соответствующая функция называется <<<P000756> с такой же подписью, как <<<P000090>.
% Эта ошибка может произойти, потому что загрузка является динамическим свойством.
Вызов функции приводит к динамической ошибке, которая возникает до
Оцениваются любые реальные аргументы.
Закрытие функции
(<<P000543>)
Это также приводит к динамической ошибке.
<<P001045>
Для каждого получателя высшего уровня <<<P000091> Оригинальное название: P000757 в
\NamespaceName{<<P001047>>>{import},i}
Соответствующий геттер по имени <<<P000758> с той же подписью, что и <<P000092>.
% Эта ошибка может произойти, потому что загрузка является динамическим свойством.
Вызов геттера приводит к динамической ошибке.
<<<P001048>
Для каждого наборщика верхнего уровня <<<P000093> Оригинальное название: P000660 в
\NamespaceName{<<P001050>>>{import},i}
Соответствующий сеттер по имени <<<P000661> с той же подписью, что и <<P000094>.
% Эта ошибка может произойти, потому что загрузка является динамическим свойством.
Вызов установщика приводит к динамической ошибке, которая возникает до
Фактический аргумент оценивается.
<<P001051>
Для каждого класса, смесителя, перечня и типа алиаса декларация под названием <<P000759> в
\NamespaceName{<<P001053>>>{import},i}
соответствующий геттер с именем <<<P000760> с возвратным типом <<P000662>.
% Эта ошибка может произойти, потому что загрузка является динамическим свойством.
Вызов геттера приводит к динамической ошибке.
<<P000510>

<<<P001054>{%>
Для того, чтобы члены импортной библиотеки могли
\NamespaceName{<<P001056>>>{отложенный}
чтобы обеспечить использование членов, которые еще не были загружены
может быть решена нормально и имеет четко определенное поведение;
Что вызывает ошибки.%
?

<<<P000801>%
При обращении к <<P000663> Успешно,
имя <<p000096> отображается в неотложенный префикс run-time namespace
\NamespaceName{\metavar{loaded}}
с привязками, как описано ниже, для немедленного импорта.
Кроме того, \NamespaceName{\metavar{loaded}}
Карты <<P000664> для функции с той же подписью, что и раньше.
Таким образом, можно ссылаться <<P000665> Опять же,
Которая всегда будет успешной.
Если вызов не удался, библиотека не была загружена.
и у вас есть возможность вызвать <<<P000666> Опять.
Неоднократный звонок по адресу <<P000667> Успехи будут разными,
как описано ниже.

<<<P001061>{%>
Обратите внимание, что это ошибка времени компиляции для отложенного префикса.
используется более чем в одном импорте;
что означает обновление связывания <<<P000100>не вмешивается в другие импортные операции.%
?

<<<P000802>%
Эффект повторного вызова <<<P000668> выглядит следующим образом:

<<P000487>
<<P001062>
При повторном вызове <<<P000669> Уже удалось,
Неоднократные призывы также увенчались успехом.
В противном случае
<<P001063>
При повторном вызове <<<P000670> Не удалось:

<<P000488>
<<P001064>
Если ошибка вызвана ошибкой компиляции,
Повторное обращение не удается по той же причине.
<<P001065>
Если неудача вызвана другими причинами,
Повторяющиеся призывы ведут себя так, как будто никакого предыдущего призыва не было.
<<P000511>
<<P000512>

<<<P001066>{%>
Другими словами, успешный <<P000671> гарантируют, что
Доступ к импорту можно получить через префикс импорта.
Если ссылка на <<P000672> Инициируется после
Успешно завершился очередной призыв,
Это также гарантирует успешное завершение (успех является идемпотентным).
Мы не указываем, какой объект возвращает будущее.%
?
<<P001067>

<<<P000803>%
<<<P001068>{Семантика немедленного импорта}
Когда <<P000104> является немедленным импортом с библиотекой URI, <<<P000105>,
в виде строки <<<P000106>:
<<<P000107> библиотека, полученная из исходного кода, обозначенного <<<P000108>.
Затем мы говорим, что URI <<<P000109> Обозначает библиотеку <<<P000110>.
Весь импорт и экспорт одного и того же URI в программе Dart обозначает
той же библиотеки,
Импорт и экспорт различных URI обозначают различные библиотеки.

<<<P000804>%
Пусть <<P001069>>>{<<P001070>>>{import}, i} будет пространством имен, импортированным из <<P000111>>>.
Пространство имен времени выполнения <<<P001071>{i} будет иметь следующие привязки:

<<P000489>
<<P001072>
Для каждой функции верхнего уровня, getter или setter <<<P000112> Название: P000113 в
\NamespaceName{<<P001074>>>{import}, i},
привязка от <<<P000114> к указанной функции, геттеру или сеттеру.
<<P001075>
Для каждого имени <<<P000761> в \NamespaceName{\metavar{import}, i}
который связан с классом, смешиванием, перечислением или объявлением псевдонима типа
Введение типа <<P000115>,
a привязка от <<<P000762> к составленному представлению <<<P000116>.
<<P000513>

<<<P000805>%
Если <<P000117> имеет приставку <<<P000118>,
Пространство имен времени выполнения текущей библиотеки
Карты <<<P000119> в неотложенное пространство имен во время выполнения \NamespaceName{p}
Содержит отображения <<<P001079>{i}.

<<<P001080>{%>
\NamespaceName{p} могут иметь дополнительные карты
Потому что может быть несколько импортов с одним и тем же префиксом.
До тех пор, пока все они являются немедленными.%
?

<<<P000806>%
В противном случае, когда <<<P000120> не имеет префикса,
Пространство имен библиотеки времени выполнения текущей библиотеки
Содержит каждое отображение в <<<P001082>{i}.
<<P001083>


<<<P001084>{Экспорт}
<<P000888>

<<<P000807>%
Библиотека <<<P000121> экспортирует пространство имен (<<P000544>), что означает, что
Заявления в пространстве имен доступны другим библиотекам.
Если вы решили импортировать <<P000122> (<<P000545>).
Пространство имен, которое <<P000123> Экспорт известен как
<<<P001085>{экспортируемое пространство имен}{библиотека!экспортируемое пространство имен}.

<<<P000808>%
Библиотека всегда экспортирует все имена и декларации в публичное пространство имен.
Кроме того, библиотека может реэкспортировать дополнительные библиотеки.
<<<P000625>, часто называемый просто <<P000626>:

<<P000490>
<libraryExport> ::= <metadata> \EXPORT{{} <configurableUri> <combinator>> "
<<P000514>

<<<P000809>%
Интерпретация настраиваемых URI описана в другом месте.
(P000546).
Экспорт определяет URI <<<P000124>
где находится декларация экспортируемой библиотеки.
Ошибка времени компиляции, если указанный URI
Это не относится к библиотечной декларации.

<<<P000810>%
The
<<P000627>
Библиотека определяется следующим образом, в два этапа.
<<<P000125> Будьте библиотекой,давать <<<P001087> E}{1}{m} являются экспортными директивами <<<P000126>,
И пусть <<<P001088>{L}{1}{m) - библиотеки
<<<P000127>> <<<P000128> для всех <<<P000129>.
Пусть \NamespaceName{\metavar{public}} будет публичным пространством имен $L$
(P000547).

<<<P001091>{%>
Обратите внимание, что приватные имена и префиксы импорта отсутствуют.
в \NamespaceName{<<P001093>>>{public}}%
?

<<<P000811>%
На первом этапе вычисляется пространство имен, предоставляемое каждой экспортируемой библиотекой:
Для каждого <<P000131>
\NamespaceName{\metavar{экспорт},i}
Пространство имен, полученное при применении
Комбинаторы пространства имен <<<P000132> то
Произошло это из-за того, что «Perfect»
(<<P000548>),
и удаление каждого связывания имени <<<P000134> такой, что
\NamespaceName{<<P001097>>>{public}} определяется по адресу <<<P000135>,
<<<P000136> и <<<P000137> Имеют одинаковое имя.
<<<P001098> Потому что местные декларации будут тенью для реэкспорта.

<<<P000812>%
The
<<<P001099>{namespace re-exported from}{library!namespace re-exported from}
<<<P000138> \NamespaceName{\metavar{экспортировано},i}.

<<<P000813>%
На втором этапе вычисляется экспортируемое пространство имен <<<P000139>.
Пусть \NamespaceName{\metavar{merged}} будет результатом применения
слияние конфликтов
(<<P000549>)
<<<P001104><<<P001105> NS}}{\metavar{экспортировано},1}{<<< P001107>>>{экспорт},m}.
Ошибка времени компиляции возникает, если имя
\NamespaceName{<<P001109>>>{merged}}
является конфликтным
(<<P000550>).

<<<P001110><%
Это правило является более строгим, чем соответствующее правило для импорта:
Когда две импортные декларации имеют столкновение названий, это всего лишь ошибка.
<<<P001111>{использовать} противоречивое название,
Не является ошибкой то, что слово «столкновение» существует.
С экспортированными именами, это ошибка, что столкновение имени существует.
Причина этого различия заключается в том, что
Конфликт может бесшумно разорвать импортеров нынешней библиотеки <<<P000140>, - говорится в сообщении.
Если мы будем использовать тот же подход к экспорту, что и к импорту:
Если библиотека <<P000141> Импорт <<<P000142>
и использует имя <<<P000143> который реэкспортируется <<<P000144> из <<<P000145>
Затем добавьте декларацию под названием <<P000146>
В реэкспортируемую библиотеку
будет использовать <<<P000147> в $L'$ Ошибка,
и хранителей <<<P000149> может быть не в состоянии изменить <<<P000150>.%
?

<<<P000814>%
The
<<<P001112>{экспортируемое пространство имен}{библиотека!экспортируемое пространство имен}
<<<P000151>> это
<<P000152>.

<<<P000815>%
Мы говорим, что имя <<P000628>
если имя находится в экспортируемом пространстве имен библиотеки.
Мы говорим, что заявление <<<P000629>
если декларация находится в экспортируемом пространстве имен библиотеки.

<<<P000816>%
Для данного <<<P000153>,
Мы говорим, что <<P000154>
<<P000630>
<<<P000155>, а также <<P000156>
<<P000631>
\NamespaceName{\metavar{экспортировано},i}.
Когда не может возникнуть путаницы, мы можем просто сказать:
<<<P000157>> <<<P001115> <<<P000158>, или
<<<P000159>> \NoIndex{re-exports} \NamespaceName{<<P001118>>>{exported},i}.


<<<P001119> Комбинаторы пространства имен
<<P000889>

<<<P000817>%
Импорт (<<P000551>>>) и экспорт (<<P000552>>>) полагаться
<<P000632>
Для настройки пространства имен
(<<P000553>)
Управляйте столкновениями имен.
Поддерживаемыми комбинаторами пространства имен являются \SHOW{} и \HIDE.

<<P000491>
<комбинатор> ::= \SHOW{} <identifierList> | \HIDE{} <identifierList>

<identifierList> ::= <identifier> (',' <identifier>) *
<<P000515>

<<<P000818>%
Мы определяем несколько операций, которые вычисляют пространства имен.
The
<<<P001124>{союз двух пространств имен}{namespace!union},
\IndexCustom{<<P000160>>>>>>{%)
$\cup$@$\NamespaceName{a} \cup \NamespaceName{b}$,
Определяется, когда каждый ключ, где оба определены, отображается на одно и то же значение.
Этот союз - пространство имен, которое отображает
Каждый ключ <<<P000163> \NamespaceName{a}
соответствующего значения \Namespace{a}{n}и каждый ключ <<<P000164> \NamespaceName{b}
<<<P001129>{b}{n}.

<<<P001130><%
Обратите внимание, что это <<<<P001131>{является} пространством имен, потому что союз только определен.
когда отображение любого заданного ключа является однозначным.%
?

<<<P000819>%
Функция
\IndexCustom{<<P001133>>>{l}{<<P001134>>>{a}}}
hide(l,NSa)@\Hide{l}{<<P001136>>>{}}}
Выберите список идентификаторов <<P000165> и пространство имен \NamespaceName{a},
и создает пространство имен \NamespaceName{b} Это то, что
идентично <<<P001139> Кроме того, что для каждого идентификатора <<<P000763> в <<<P000166>,
<<<P001140>{b} не определено в <<<P000764> <<<P000673>.

<<<P000820>%
Функция
\IndexCustom{<<P001142>>>{l}{<<P001143>>>{a}}}
show(l,NSa)@\Show{l}{<<P001145>>>{}}}
Возьмите список идентификаторов <<P000167> и пространство имен <<<P001146>{a},
и создает пространство имен <<<P001147> это
Карты каждого идентификатора <<<P000765> <<<P000168>
где \NamespaceName{a}
то <<<P001149>{a}{n}.
Кроме того, для каждого идентификатора <<<P000766> <<<P000169>
где \NamespaceName{a} определено в <<<P000674>,
\NamespaceName{b} карты \code{\id=} до \Namespace{a}{<<P000676>>>}.
Наконец, <<<P001153>{b} не определено во всех других именах.

<<<P000821>%
Пусть \List{C}{1}{n} будет последовательностью комбинаторных оговорок.
<<<P001155> Они будут поступать из импортной или экспортной директивы. ?
Результатом
<<<P001156>{применение комбинаторных оговорок}{%
Комбинаторные предложения! Приложение для пространства имен
в заданное пространство имен \NamespaceName{\metavar{start}}
является пространством имен \NamespaceName{\metavar{end}},
который рассчитывается следующим образом.

<<<P000822>%
Пусть \NamespaceName{0} будет \NamespaceName{\metavar{старт}}.
Для каждого комбинатора <<<P000170>, <<P000171>:
Если <<P000172> имеет форму \code{\SHOW{} <<P000173> Тогда пусть
<<P000174>.
Если <<P000175> имеет форму \code{\HIDE{} <<P000176> Тогда пусть
<<P000177>.
Тогда \NamespaceName{\metavar{end}} является \NamespaceName{n}.

<<<P001171>{%>
\NamespaceName{\metavar{end}} всегда соглашается с
\NamespaceName{\metavar{старт}} везде, где оба определены,
и набор имен, где определен \NamespaceName{\metavar{end}}
всегда является подмножеством того из <<<P001178>{<<<P001179>{старт}}.
В этом смысле это процедура <<<P001180>{сужение}.%
?

<<<P001181>{%>
Обратите внимание, что можно использовать \HIDE{} или <<P001183>>>{} на идентификаторе
которого нет в названном пространстве.%
?
<<<P001184>{%>
Это позволяет избежать ситуаций, когда, например,
удаление декларации из библиотеки <<<P000178> вызвать
Библиотека, которая импортирует <<<P000179>%
?


<<<P001185> Слияние пространств имен.
<<P000890>

<<<P000823>%
В этом разделе мы определяем операцию на пространствах имен.
которые устраняют определенные привязки имен в случае столкновений имен,
и отслеживает такие столкновения имен, используя особое значение,
<<P000633>.

<<<P000824>%
Когда имя <<P000180> отображается на <<<P001186>{} по пространству имен,
Мы говорим: <<P000181> это
\IndexCustom{conflicted}{конфликтное название}.
Поиск противоречивого имени будет успешным
Нормальные правила для пространств имен и лексического определения,
Но это ошибка времени компиляции в месте, где используется имя.
Это название является противоречивым.

<<<P000825>%
Пусть \List{<<P001189>>>{NS}}{1}{m} будет списком пространств имен
<<<P001190>{(список может содержать или не содержать дубликатов)}.
The
<<<P001191>{конфликтное слияние}{именное пространство!конфликтное слияние}
<<<P001192><<<P001193> NS}}{1}{m} является
$\NamespaceName{<<P001195>>>{конфликт}}\cup
\NamespaceName{\metavar{сокращенно},1} <<<P001199><<<P001728>
\ldots\ \NamespaceName{\metavar{narrowed},m}$,
где пространство имен \NamespaceName{<<P001204>>>{конфликт}} и
\NamespaceName{\metavar{narrowed},i}, <<P000182>,
Указывается следующее:

<<<P000826>%\NamespaceName{\metavar{конфликт}} пустой, за исключением
Обязательства, упомянутые ниже.
Для каждого <<<P000183>,
\NamespaceName{\metavar{суженный},i} идентичен \NamespaceName{i},
За исключением того, что первое не определено в каждом имени <<<P000184> где
Последняя определяется,
и одно из следующих условий:

<<P000492>
<<<p001212> Существует <<<P000185> Название: P000186 такой, что
<<<P001213>{j} определено в <<<P000187>,
<<P000188> и <<P000189> иметь одинаковые базовые имена,
<<<P001214>{i}{n} — декларация в системной библиотеке,
и <<<P001215>{j}{n'} является декларацией в несистемной библиотеке.
<<<<P001216>{%>
Заявление из несистемной библиотеки теней
Заявления из системных библиотек.%
?
<<P001217> В противном случае существует <<<P000190> Название: P000191 такой, что
\NamespaceName{j} определено в <<<P000192>,
<<<P000193> и <<<P000194> иметь одинаковые базовые имена,
\Namespace{i}{n} и <<P001220>>>{j}{n'} не являются одним и тем же заявлением
И не одноимённое пространство
а не геттер и сеттер, объявленные в той же библиотеке,
Ни один из них, ни оба
\Namespace{i}{n} и \Namespace{j}{n''
Деклараций в системной библиотеке.
%
<<<P001223>{%>
С двумя разными объявлениями с одинаковым именем,
Оба они устранены,
За исключением случаев, когда это геттер и сеттер из одной и той же библиотеки.
Обратите внимание, что \Namespace{i}{n} и \Namespace{j}{n'')
оба должны быть объявлениями или оба должны быть пространствами имен;
потому что слияние конфликтов применяется к пространствам имен, где
импортные декларации, противоречащие импортному префиксу
Они уже ликвидированы.
Когда они оба пространства имен, нет никакого конфликта.
Потому что это одно и то же имя.%
?

В этой ситуации <<<P000195> Соответствует <<<P001226>
\NamespaceName{<<P001228>>>{конфликт}}
<<<P001229>{%>
Мы можем поменять все имена и снова использовать правило.
<<<P000196> Конфликт.%
?
<<P000516>

<<<P001230>{%>
Короче говоря, слияние конфликтов берет список пространств имен и производит
Пространство имен, которое является объединением нескольких несвязанных пространств имен
(Поскольку все названия были уничтожены):
Пространство имен конфликтов, которое записывает все неразрешенные столкновения имен.
и список пространств имен, где были удалены сталкивающиеся имена.
Некоторые столкновения названий были решены путем предпочтения.
Заявления из несистемных библиотек
по заявлениям системных библиотек.%
?

<<<P000827>%
Полезно уметь ссылаться на результат сужения.
Пусть \List{<<P001232>>>{NS}}{1}{m} будет списком пространств имен, $i \in 1 .. m$,
и рассматривать слияние конфликта, как указано выше.
The
<<<P001233>{конфликт сузил пространство имен}{namespace!conflict narrowed}
Из <<<P001234>{i} получается \NamespaceName{<<<P001236>{narrowed},i}.


<<<P001237> Части.
<<P000891>

<<<P000828>%
Библиотеку можно разделить на <<<P000634>,
Каждый из них может храниться в отдельном месте.
Библиотека идентифицирует свои части, перечисляя их через <<<P001238>{} директивы.

<<<P000829>%
<<<P000635> Укажите URI, где
а Блок компиляции Dart, который должен быть включен в текущую библиотеку
Может быть найден.

<<P000493>
<partDirective> ::= <metadata> <<<P001239>{{i}>> "

<part Header> ::= <metadata> <<<P001240>{} \OF{} (<dottedIdentifierList> | <uri>); "

<parteclaration> ::=
<partHeader> (<metadata> <topLevelDeclaration>)* <EOF>
<<P000517>

<<<P000830>%
Часть содержит строку, полученную из <<<P001242>{partDeclaration}.

<<<P001243>{%>
Таким образом, можно сказать, что <<<P001244>{partDeclaration} является начальным символом грамматики.
Как описано в статье.~\ref{librariesAndScripts}.%
?

<<<P000831>%
<<<P000636> начинается с <<<P001245> <<<P001246>< сопровождаемый
Название библиотеки, к которой принадлежит часть,
или <<<P001247>{uri}, обозначающий указанную библиотеку.
Декларация части состоит из заголовка части, за которым следуетПоследовательность деклараций высшего уровня.

<<<P000832>%
Составление части директивы формы \code{\PART{) <<<P000198>;
Это приводит к тому, что система Дарта пытается собрать содержимое
URI - значение <<<P000199>.
Объявления верхнего уровня на этом URI затем компилируются компилятором Dart.
в рамках действующей библиотеки.
Ошибка времени компиляции, если содержимое URI не
декларация действительной части.
Это ошибка времени компиляции, если указанная часть декларации <<<P000200> имена
библиотека, отличная от текущей библиотеки, в которой <<<P000201> принадлежит.

<<P000833>%
Ошибка времени компиляции, если библиотека содержит
Две части директивы с одинаковым URI.

<<<P000834>%
Мы говорим, что библиотека <<<P000202> является <<<P000637> Библиотека <<<P000203> если
все нижеследующее верно (<<<P000555>, <<P000556>:

<<P000494>
<<P001250> <<<P000204> и <<P000205> — одна и та же библиотека.
<<P001251> <<P000206> импортирует или экспортирует библиотеку <<<P000207> и <<<P000208> Доступно из <<<P000209>.
<<P000518>

<<<P000835>%
<<<P000210> Будьте библиотекой, пусть <<<P000211> Будь Ури,
<<<P000212> и <<<P000213> Отдельные библиотеки, доступные из <<<P000214>.
Ошибка времени компиляции, если <<<P000215> и <<<P000216> Оба содержат
директива с URI <<<P000217>.

<<<P001252>{%>
В частности, ошибкой является использование одной и той же части дважды.
В той же программе
(<<P000557>).
Обратите внимание, что относительный URI интерпретируется как относительный по отношению к местоположению URI.
библиотека (<<P000558>), что означает <<P000218> и <<P000219> могут оба
имеют часть, идентифицированную по <<<P000677>, но они не одинаковы
URI, если <<<P000220> и <<<P000221> имеют одинаковое месторасположение.%
?


<<<P001253>{Сценарии}
<<P000892>

<<P000836>%
<<<P000638> Библиотека, экспортировавшая пространство имен (<<P000559>) включает
Декларация функций верхнего уровня <<<P000678>
Для этого требуется либо ноль, либо два аргумента.

Сценарий <<<P000222> Выполняется следующим образом:

<<<P000837>%
Во-первых, <<<P000223> составляется как библиотека, как указано выше.
Функция верхнего уровня, определяемая как <<P000679>
в экспортируемом пространстве имен <<<P000224> (перенаправлено с «P000560»)
следующим образом:
Если <<P000680> может быть вызвана двумя позиционными аргументами,
Он ссылается на следующие два фактических аргумента:

<<P000495>
<<<P001254> Объект, тип времени выполнения которого реализует <<P000681>.
<<P001255> Объект, указанный в текущем Изоляторе <<<P000225> был создан,
Например, посредством ссылки <<<P000682> что породило <<<P000226>,
или нулевого объекта (<<<P000561>), если такой объект не был поставлен.
<<P000519>

<<<P000838>%
Если <<P000683> не может быть вызван двумя аргументами,
Но его можно назвать одним позиционным аргументом.
он вызывается объектом, чей тип времени выполнения реализует <<P000684>
Как единственный аргумент.
Если <<P000685> не может быть вызван одним или двумя аргументами,
На него ссылаются без аргументов.

<<<P001256>{%>
Обратите внимание, что если <<<P000686> требует более двух аргументов,
Библиотека не считается скриптом.%
?

<<<P001257>{%>
Программа Dart обычно выполняется путем выполнения сценария.%
?

<<<P000839>%
Ошибка времени компиляции, если объем экспорта библиотеки содержит декларацию
Называется она «P000687», а библиотека — не сценарий.
<<<P001258>{%>
Это ограничение гарантирует, что все топ-уровни <<P000688> заявления
Введите главную функцию сценария, чтобы не было
Геттер верхнего уровня или поле
не может быть функцией, требующей более двух
аргументы. Ограничение позволяет инструментам рано выходить из строя на недействительном <<<P000690>
без необходимости знать, будет ли библиотека использоваться в качестве входа
Смысл программы Дарт. Возможно, это ограничение будет снято.
в будущем.%
?<<<P001259> УРИ
<<P000893>

<<<P000840>%
URI определяются с помощью струнных букв:

<<P000496>
<uri>:= <stringLiteral>

<configurableUri> ::= <uri> <конфигурация Ури*

<configurationUri> ::= \IF{} '(' <uriTest> ''

<uriTest> ::= <dottedIdentifierList> ('==' <stringLiteral>)?
<<P000520>

<<<P000841>%
Это ошибка времени компиляции, если строка буквально описывает URI.
или строка, которая используется в <<<P001261>
содержит интерполяцию строк.

<<<P000842>%
<<<P000639> <<<P000227> в форме
\code{<<P001263>>>{uri) <<P000228> <<P001264> <<P000229>
<<<P001265>{указывает URI}{указывает URI} следующим образом:

<<P000497>
<<P001266>
Пусть <<<P000230> будет \metavar{uri}.
<<P001268>
Для каждой из следующих конфигураций URI формы
\code{<<P001270>>>{) (<<P000231>) <<P000232>,
В порядке источника сделайте следующее.

<<P000498>
<<P001271>
Если <<<<P000233> \code{\metavar{ids}}
без <<<P001274> Оговорка, это
эквивалентно \code{\metavar{ids} == "истинное"}.
<<<P001277>
Если $\metavar{test}_i$ \code{\metavar{ids} == \metavar{string}}
Затем создайте строку \metavar{key} из \metavar{ids}
путем конкатенирования идентификаторов и точек,
Опустить любые пробелы между ними, которые могут возникнуть в источнике.
<<<P001283>
Смотрите <<<P001284>{key} в доступном разделе
<<P000640>.
<<<P001285>{%>
Среда компиляции предоставляется платформой.
Он отображает некоторые струнные ключи для значений строк,
могут быть получены программным путем с использованием
<<P000691> Конструктор.
Инструменты могут выбирать только некоторые части среды компиляции.
Доступен для выбора конфигурации URIs.%
?
<<P001286>
Если окружение содержит запись для <<<P001287>{ключ} и
ассоциированная стоимость равна, как постоянная струнная стоимость, значению
строка буквальная \metavar{string},
<<<P000235> <<<P000236> Перестаньте повторять URI конфигурации.
<<<P001289>
В противном случае переходите к следующей конфигурации URI.
<<P000521>
<<P001290>
URI, указанный в <<<P000237>, является <<<P000238>.
<<P000522>

<<<P000843>%
Эта спецификация не обсуждает интерпретацию URI.
со следующими исключениями.

<<<P001291>{%>
Интерпретация URI в основном оставлена на
окружающей вычислительной среды.
Например, если Dart работает в веб-браузере,
Этот браузер, скорее всего, интерпретирует некоторые URI.
Хотя может показаться привлекательным указать, скажем, что
URI интерпретируются в отношении стандарта, такого как IETF RFC 3986.
На практике это обычно зависит от браузера.
На него нельзя положиться.%
?

<<<P000844>%
URI формы <<<P000692> толкуется как
Ссылка на системную библиотеку (<<P000562>) <<P000240>.

<<<P000845>%
URI формы <<<P000693> интерпретируется в
конкретного способа реализации.

<<<P001292>{%>
Цель состоит в том, чтобы во время разработки программисты Dart могли полагаться на
менеджер пакетов для поиска элементов своей программы.%
?

<<<P000846>%
В противном случае любой относительный URI интерпретируется как относительный.
Местонахождение нынешней библиотеки.
Вся дальнейшая интерпретация URI зависит от реализации.

<<<p001293>> Это означает, что он зависит от времени выполнения.


<<<P001294> Типы
<<P000894>

<<<P000847>%
Dart поддерживает статическую типизацию на основе типов интерфейсов.

<<<P001295>{%>
Типовая система является звуковой в том смысле, что
если переменная типа <<<P000242> относится к объекту типа <<<P000243> во время бега,
<<<P000244> Это подтип <<<P000245>.
Другими словами, содержимое кучи удовлетворяет ожиданиям.
выражается статической типизацией.

Однако параметры типа являются ковариантными.
(например, \SubtypeNE{\code{List<int>}}{\code{List<num>}})Это означает, что некоторые операции подвергаются динамическим проверкам типа.
(например, <<<P000696>, который будет бросаться во время бега)
Если <<P000697> объявленный тип <<P000698>,
На самом деле это P000699.
Таким образом, программа может быть свободна от ошибок времени компиляции.
И он все еще может повлечь за собой ошибку типа во время выполнения.

Этот выбор был сделан намеренно в первые дни Дарта.
(И это, несомненно, спорно).
Представляет собой компромисс, где
Потенциал для ошибок типа времени выполнения - это стоимость.
Преимуществом являются более простые программы.
В последние годы Dart развивался, чтобы иметь все больше и больше безопасности статического типа.
Например, нулевая безопасность, и эта тенденция, вероятно, продолжится.
?


<<<P001297> Статические типы
<<P000895>

<<<P000848>%
Аннотации типов могут встречаться в переменных декларациях (<<P000563>),
включая формальные параметры (<<P000564>),
в возвратных типах функций (<<<P000565>),
и в границах переменных типа (<<P000566>).
Аннотации типов используются во время статической проверки и при запуске программ.
Типы определяются с использованием следующих правил грамматики.

<<<P001298>{%>
<<<P001299>{typeIdentifier} - идентификатор, который может быть именем типа,
То есть,
обозначает <<<P001300> Идентификатор, который не является <<<P001301> BUILT<<<P001730> В <<<P001731> Идентификатор
(<<P000567>)%
?

<<<P001302>{%>
Нетерминалы с названиями формы \synt{\ldots{}NotFunction}
Выведите термины, которые являются типами, которые не являются типами функций.
Обратите внимание, что в нем <<<P001305>{does} выводится тип <<<P001306>,
который сам по себе не является функцией,
Но это наименьшая верхняя граница всех типов функций.
?

<<P000499>

<type> ::= <functionType> '?'?
<<<P001307> <typeNotFunction>

<typeNotVoid> ::= <functionType> '?'
<<<P001308> <typeNotVoidNotFunction> '?'

<typeNotFunction> ::= <<<P001309><
<<<P001310> <typeNotVoidNotFunction>>?

<typeNotVoidNotFunction> ::= <typeName> <typeArguments>
<<<P001311> (<typeIdentifier>>)? \FUNCTION{}

<typeName> ::= <typeIdentifier> ('.' <typeIdentifier>)?

<typeArguments> ::= <<typeList>> "

<typeList> ::= <type> (',' <type>)*

<typeNotVoidNotFunctionList>:= <<<P001313><
<typeNotVoidNotFunction> (',' <typeNotVoidNotFunction>) *

<functionType> ::= <functionTypeTails>
<<<P001314> <typeNotFunction> <functionTypeTails>

<функция TypeTails>::= <functionTypeTail> '?'? <functionTypeTails>
<<P001315> <functionTypeTail>

<функция TypeTail ::= \FUNCTION{} <typeParameters> <parameterTypeList>

<parameterTypeList> ::= ''' "
<<<P001317><<normalParameterTypes>> <optionalParameterTypes>> "
<<<P001318> <<типы нормального параметра>> "
<<<P001319> <("Необязательные параметры">> "

<normalParameterTypes> ::= <<<P001320><
<normalParameterType> (>normalParameterType>)*

<normalParameterType> ::= <metadata> <typedIdentifier
<<<P001321> <metadata> <type>

<optionalParameterTypes> ::= <optionalPositionalParameterTypes>
<<<P001322> <namedParameterTypes>

<optionalPositionalParameterTypes> ::= '[' <normalParameterTypes> ','?' "

<namedParameterTypes> ::=
\{> <namedParameterType> ('', <namedParameterType>''), \} "

<namedParameterType> ::=
<metadata> <<P001323> <typedIdentifier>

<typedIdentifier> ::= <type> <identifier>
<<P000523>

<<<P000849>%
А. Реализация Dart должна обеспечивать статическую проверку, которая обнаруживает и сообщает
именно те ситуации, которые эта спецификация идентифицирует как ошибки времени компиляции;
И только эти ситуации.
Аналогично, статическая шашка должна выдавать статические предупреждения.
По крайней мере, для ситуаций, указанных как таковые в данной спецификации.

<<<P001324>{%>
Ничто не исключает дополнительных инструментов, которые реализуют альтернативный статический анализ.
(например, толкование аннотаций существующего типа в звуковой форме)
такие как невариантные дженерики или вывод о дисперсии на основе декларации
из фактических деклараций.Однако использование этих инструментов не должно препятствовать
Успешная компиляция и выполнение кода Dart.%
?

<<<P000850>%
Тип <<<P000246> <<<P000641> Если [f]:

<<P000500>
<<P001325>
<<P000247> Имеет форму <<P000767> или форму \code{\metavar{префикс}.\id},
И это не означает объявление типа.
<<P001329>
<<P000248> обозначает переменную типа,
Это происходит в сигнатуре или теле статического элемента.
<<P001330>
<<P000249> является параметризованным типом формы <<<P000700>,
<<P000252> является неполноценным,
<<<P000253> не является общим типом,
<<<P000254> является родовым типом,
Об этом сообщает <<P000255> Типовые параметры и <<P000256>,
<<<P000257> Для некоторых <<<P000258>.
<<P001331>
<<P000259> является функциональным типом формы

<<<P001332><<<P000260> \FUNCTION{}<$X_1\ \EXTENDS\ B_1, \ldots,\ X_m\ \EXTENDS\ B_m$>

<<P000701>

<<P001334>
или в форме

<<<P001335><<<P000264> \FUNCTION{}$X_1\ \EXTENDS\ B_1, \ldots,\ X_m\ \EXTENDS\ B_m$>

<<<P001337><<<P001338> $T_1\ x_1, \ldots,\ T_k\ x_k,\ $\{$T_{k+1}\ x_{k+1}, \ldots,\ T_n\ x_n$\}]

<<P001339>
где каждый <<P000268> который не является поименованным параметром, может быть опущен;
<<P000269> является деформированным для некоторых <<<P000270>,
<<<P000271> Для некоторых из них <<P000272>.
<<P001340>
<<P000273> Деклараций, которые были импортированы из нескольких пунктов импорта.
<<P000524>

<<<P000851>%
Любое появление неправильного типа в библиотеке является ошибкой времени компиляции.

<<<P000852>%
Тип <<<P000274> \IndexCustom{deferred}{type!deferred}
Если он имеет форму <<<P000275> <<<P000276> Это отложенный префикс.
Ошибка времени компиляции для использования отложенного типа
в аннотации типа, тесте типа, отливке типа или в качестве параметра типа.
Тем не менее, все другие ошибки времени компиляции должны быть выпущены.
Все отложенные библиотеки были успешно загружены.

% Теперь, когда передается дженерик, p.T также должен рассматриваться как динамический -
В противном случае мы должны немедленно потерпеть неудачу. Где мы это скажем? И как это делается
% это согласуется с идеей о том, что как объект типа он терпит неудачу? Стоит ли говорить, что
%-ный аксессуар на p возвращает динамику вместо отказа? Различают ли их использование
% в конструкторе против его использования в аннотации? Не то чтобы мы оценивали тип
% объектов в конструкторских аргах — они не могут представлять параметризованные типы.


<<P001342> Тип поощрения
<<P000896>

<<<P000853>%
Система статического типа приписывает статический тип каждому выражению.
В некоторых случаях тип локальной переменной
<<<P001343>{(который может быть формальным параметром)}
могут быть повышены из заявленного типа, исходя из управляющего потока.

<<<P000854>%
Мы говорим, что переменная <<<P000277> Известно, что тип <<P000278>
Каждый раз, когда мы разрешаем рекламировать тип <<<P000279>.
Точные обстоятельства, когда разрешено продвижение по типу, приведены в
соответствующие разделы спецификации
(<<P000568>, <<P000569> и <<P000570>.

<<<P000855>%
Продвижение типа для переменной <<<P000280> Допускается только тогда, когда мы можем сделать вывод, что
Такое продвижение действительно на основе анализа определенных булевых выражений.
В таких случаях мы говорим, что
Булево выражение <<P000281> Показывает, что <<<P000282> имеет тип <<P000283>.
Как правило, для всех переменных <<<P000284> и типы <<<P000285>, булево выражение
Не показывает, что <<<P000286> имеет тип <<P000287>.
Те ситуации, когда выражение действительно показывает, что переменная имеет тип
прямо упомянутых в соответствующих разделах настоящей спецификации
(<<P000571> и <<P000572>).


<<P001344> Система динамического типа
<<P000897>

% <<P000573> Когда имя класса появляется как тип, %
% это имя обозначает интерфейс класса. Так что мы можем сказать
% что динамический тип экземпляра является негенерическим классом или
% генерическая инстанциация генерического класса.

<<<P000856>%
Пусть <<<P000288> будет примером.<<<P000642> <<<P000289> Класс, который определяется
Для ситуации, когда <<P000290> был получен в качестве нового примера
(<<P000574>,
<<P000575>, <<P000576>, <<P000577>, <<P000578>.

<<<P001345>{%>
В частности, динамический тип экземпляра никогда не меняется.
Иногда указывается только, что данный класс реализует
определенный тип, например, для списка буквальный.
В этих случаях динамический тип зависит от реализации.
За исключением, конечно, того, что упомянутое ограничение суперинтерфейса должно быть удовлетворено.
?

<<<P000857>%
Динамические типы запущенной программы Дарта эквивалентны
Статические типы в отношении субтипирования.

<<<P001346>{%>
Определенные динамические проверки типа выполняются во время выполнения.
(<<P000579>,
<<P000580>,
<<P000581>,
<<P000582>,
<<P000583>,
<<P000584>,
<<P000585>,
<<P000586>,
<<P000587>,
<<P000588>,
<<P000589>,
<<P000590>,
<<P000591>,
<<P000592>,
<<P000593>,
<<P000594>.
Как указано в этих местах,
Эти динамические проверки основаны на динамических типах экземпляров.
и фактические виды деклараций
(<<P000595>)%
?

<<<P000858>%
Когда типы уточнены как экземпляры встроенного класса <<<P000702>,
Эти объекты перекрывают <<<P001347> оператор
Унаследовано от <<<P000703> класс, так что
<<<P000704> Объекты равны по оператору \lit{=======================================================================================================================================================================================================================================================
f соответствующие типы являются подтипами друг друга.

<<<P001349>{%>
Например, <<<P000705> объекты для типов
\DYNAMIC{{} и <<P000706>>> равны между собой.
<<<P000707>> Вы должны оценить до <<<P001351>.
На встроенную функцию не накладываются ограничения <<<P000708>,
<<<P000709> может быть <<<P001352>{} или <<<P001353>.

Точно так же <<P000710> Примеры заявлений о различных типах псевдонимов
Объявление имени для одного и того же типа функций равно: %
?

<<P000000>

<<<P000859>%
<<<P001354>{%>
Случаи <<<P000711> могут быть получены различными способами,
Например, используя рефлексию,
Читать далее «P000712» Для объекта,
или путем оценки выражения <<<P001355>.%
?

<<<P000860>%
Выражение представляет собой <<<P001356> [тип буквальный], если оно является идентификатором,
квалифицированный идентификатор,
который обозначает класс, смесь, enum или тип alias Declaration, или это
идентификатор, обозначающий параметр типа общего класса или функции.
Это <<<P001357>{постоянный тип буквальный} если он не обозначает параметр типа,
и не квалифицируется отложенным префиксом.
<<<P001358>{%>
Постоянная буквальная форма - это постоянное выражение (<<P000596>).%
?


<<P001359> Типовые алиазы
<<P000898>

<<<P000861>%
<<<P000643> Назовите имя какого-нибудь типа,
или для отображения от аргументов типа к типам.

<<<P001360>{%>
Для такого заявления обычно используется фраза "a typedef",
Из-за заметного появления токена <<<P001361>>%
?

<<P000501>
<typeAlias> ::= \gnewline{}
<<<P001363>{} <typeIdentifier> <typeParameters>? '=' <type>'; "
<<P001364> <<<P001365>< <functionTypeAlias>

<functionTypeAlias> ::= <functionPrefix> <formalParameterPart>; "

<functionPrefix> ::= <type>? <identifier>
<<P000525>

<<<P000862>%
Рассмотрим заявление типа alias <<P000291> в форме

<<P001366>
\code{<<P001368>>>{) \id\TypeParametersStd> = $T$;

<<P001371>
Объявлено в библиотеке <<P000293>.
Эффект <<<P000294> Предлагается ввести <<<P000768> в библиотечный объем <<<P000295>.
Когда <<P000296>
<<<P001372>{(где тип псевдонима не является общим)}
<<P000769> Связано с <<<P000297>.
Когда <<P000298>
\commentary{(где тип псевдонима является общим)}
<<P000770> привязан к отображению из списка аргументов типа
<<<P001374>{U}{1}{}
к типу
<<P000299>.

<<<P000863>%
В предположении, что <<<P001375>{X}{1}{s} являются такими типами, что<<<P000300>, для всех <<<P000301>,
Ошибка времени компиляции <<<P000302> не является регулярным,
и это ошибка времени компиляции, если какой-либо тип встречается в <<<P000303> Он не очень хорошо ограничен.

<<<P001376>{%>
Это означает, что границы, объявленные для
Формальные параметры типа общего псевдонима типа
должны быть такими, чтобы, когда они удовлетворены,
границы, относящиеся к телу, также удовлетворены;
и тип, происходящий как субтермин тела, может нарушать его границы;
но только в том случае, если это правильный сверхограниченный тип.
?

<<<P000864>%
Кроме того,
<<<P000304> быть типами
<<<P000305>> быть параметризованным типом <<P000713>
в месте, где <<P000771> обозначает <<P000307>.
Ошибка времени компиляции <<<P000308>.
Ошибка времени компиляции <<<P000309> не является хорошо ограниченным
(<<P000597>).

<<<P000865>%
По историческим причинам тип псевдонима может иметь еще две формы.
производное использование <<<P001377>{функцияТипАлиас}.
Пусть <<<P001378>{S?} будет термином, который является пустым или полученным из <<<P001379>{тип},
и аналогично для <<<P001380>{T_j?} для любого <<<P000310>.
Затем две более старые формы определяются следующим образом:

<<<P000866>%
Типовой псевдоним формы

<<P001381>
\code{<<P001383>>>{) <<P000311> \id<\TypeParametersStd>()

<<P001386>
<<P000714>

<<P001387>
рассматривается как

<<P001388>
\code{<<P001390>>>{) <<<P001391><<<<P001392>>

<<< 001393>
<<P000715>

<<<P000867>%
Типовой псевдоним формы

<<P001394>
\code{<<P001396>>>{) <<P000315> <<<P001397><\TypeParametersStd>()

<<P001399>
<<P000716>

<<P001400>
рассматривается как

<<<P001401>
\code{<<P001403>>>{} <<<P001404><<<<P001405>>

<<<P001406>
<<P000717>

<<<P000868>%
В этих правилах для каждого $j$, если $T_j?$ пустой
затем <<<P001407>{T_j} - <<<P001408>,
иначе <<<P000321> <<<P000322>.

<<<P001409>{%>
Это означает, что старые формы позволяют опустить тип параметра.
В этом случае оно считается <<<P001410>,
Но имя параметра не может быть опущено.
Такое поведение подвержено ошибкам,
и, следовательно, рекомендуется более новая форма (с <<<P001411>{=}).
Например, заявление \code{\TYPEDEF{} <<<P000323>(int);
Уточняется, что <<P000324> обозначает тип
\code{<<P001415>>>{) \FUNCTION(\DYNAMIC)
документируется, но технически игнорируется,
Этот параметр имеет имя <<P000718>.
Очень вероятно, что читатель неправильно это прочитает.
Предположим, что <<P000719> Тип параметра.%
?

<<<P000869>%
Это ошибка времени компиляции, если значение по умолчанию указано для
формальный параметр в этих старых формах,
или если формальный параметр имеет модификатор <<<P001418>.

<<<<P001419>{%>
Обратите внимание, что старые формы могут определять только типы функций.
Они не могут определить тип общей функции.
Когда такой тип псевдонима имеет параметры типа,
Она всегда выражает семейство негенерических типов функций.
Эти ограничения существуют, потому что синтаксис был определен.
До добавления общих функций в Dart.%
?

<<<P000870>%
Пусть <<<P000325> быть типом alias заявление формы

<<<P001420>
\code{<<P001422>>>{) $F$<\TypeParametersStd> = $T$;

<<P001424>
Если <<P000328> или <<P000329> Для некоторых <<P000330> является или содержит
а \synt{typeName}} $G$ обозначение заявления типа alias <<P000332>,
Мы говорим, что
<<<P001426><<<P000333> Зависит от <<<P000334>> [тип псевдонима]! зависимость.

%%TODO(Eernst): Расслабление, когда возникает зависимость типа псевдонима
% % через <<P000335>, но не через <<<P000336>, является более разрешительным и позволяет
% % тип псевдонимов, которые должны быть объявлены, что будет ошибкой с вышеуказанным правилом.
% % Очень вероятно, что можно использовать более расслабленное правило. Это а
% возможного улучшения в будущем и непреодолимых изменений.

<<<P000871>%
<<<P000337> быть типом псевдонима декларации,И пусть <<P000338> Промежуточное закрытие
Заявления типа alias, которые <<<P000339> Зависит от.
Ошибка времени компиляции возникает, если <<P000340>.

<<<P001427>{%>
Другими словами, это ошибка для объявления псевдонима типа.
зависеть от себя, прямо или косвенно.%
?

<<<P001428>{%>
Эта ошибка может возникнуть, когда аргументы типа
Опущены в программе, но добавляются при статическом анализе
Через инстанциацию связываться
(<<P000598>)
или с помощью типоразмера,
который будет указан позднее
(<<P000599>)%
?

<<P000872>%
Когда <<P000341> является типом alias декларации формы

<<P001429>
\code{\TYPEDEF{} $F$<\TypeParametersStd> = $T$;

<<P001433>
Параметризованный тип <<<P000344> в форме
<<<P001434><<<P000345><<<<P001435> U}{1}{}>}
в области, где <<<P000346> Выберите <<P000347>
<<<P001436>{alias расширяется в один шаг до}{%
Типовой псевдоним! alias расширяется в один шаг
<<P000348>.

<<<P001437>{%>
Обратите внимание, что <<P000349> может быть равно нулю, в этом случае <<P000350> не является общим,
и мы просто заменяем имя псевдонима на его тело.
?

<<P000873>%
Если <<P000351> Это тип, который мы можем многократно заменять каждый субтермин <<P000352>
который является параметризованным типом, применяющим псевдоним типа к некоторым аргументам типа
Расширение в один шаг
<<<P001438>{(в том числе необычный случай, когда нет аргументов типа)}.
Когда дальнейшие шаги невозможны, мы говорим, что полученный тип
<<<P001439>{расширение транзитивного псевдонима}{%
Типовой псевдоним! Промежуточное расширение псевдонима
<<P000353>.

<<<P001440>{%>
Обратите внимание, что транзитивное расширение псевдонима существует,
потому что это ошибка для объявления псевдонима типа
Зависить от самого себя.
Однако он может быть большим.
Например, \code{\TYPEDEF{} <<<P000354><<<P000355>> = Карта<<<P000356> <<P000357>> доходность
увеличение экспоненциального размера для \code{$F$<$F$<\ldots$F$<<\ldots{}>{}>>.%
?

<<<P000874>%
<<<P000902>%
<<<P000361> быть типом alias заявление формы

<<P001446>
\code{<<P001448>>>{) <<<P000362><<<<P001449>> = <<<P000363>;

<<<P001450>
И пусть <<<P001451> U - тип формы <<P000364> или <<P000720>
в области, где этот термин обозначает <<P000367>.
Предположим, что переходное псевдонимное расширение \code{$F$<\List{ X}{1}{s}>} это
<<<P000903>%
Тип одной из форм <<P000369> или <<P000721>,
необязательно с последующим \synt{typeArguments},
где <<P000372> является идентификатором, обозначающим префикс импорта,
<<P000373> <<<P000722>> обозначает класс или смесь
<<<P001455>{(в частности, <<P000376> не может быть переменной типа.
Предположим, что <<P000377> встречается в выражении <<P000378> в форме
<<<P001456><<<P000379>. \id\;\metavar{args}} "
где \metavar{args} происходит от \syntax{<argumentPart>?},
<<<P000772> Это имя
Статический член <<<P000380> соответственно <<P000723>.
Выражение <<P000383> Затем рассматривается как
<<<P001461><<<P000384>. \id\;\metavar{args}} "
соответственно
<<<P001464><<<P000385>. <<P000386>. \id\;\metavar{args}}.

<<<P001467>{%>
Это означает, что можно использовать псевдоним типа
Для вызова статического члена класса или смесителя.
Например: %
?

<<<P000001>

<<<P001468>{%>
Обратите внимание, что аргументы типа перешли на <<<P000387> соответственно <<P000724>
стираются,
Таким образом, полученное выражение может быть правильным статическим вызовом члена.
Если будущая версия Dart позволяет передавать аргументы типа через класс
при обращении к статичному члену,
Предпочтительнее сохранить эти аргументы.
В то время,
существующие призывы, когда аргументы типа передаются таким образомНе сломается,
потому что в настоящее время для статического элемента является ошибкой зависеть от значения
формальный параметр типа закрывающего класса.%
?

<<P000875>%
<<<P000904>%
<<<P000390> быть типом alias заявление формы

<<P001469>
\code{<<P001471>>>{) $F$<\TypeParametersStd> = $T$;

<<P001473>
где <<P000393>.
<<<P000905>%
1990-1474-й год. Y}{1}{s} - переменные свежего типа,
Предполагалось, что он удовлетворяет пределам <<<P000394>.
Мы говорим, что <<P000395>
<<<P001475>{расширяется до переменной типа}
Если транзитивный псевдоним расширение
<<<P001476><<<P000396><<<<P001477> Y}{1}{}>}
<<<P000397> для некоторых <<<P000398>.

<<<P000876>%
<<<P000906>%
<<<P000399> быть параметризованным <<<P001478><<<P000400><<<<P001479> T}{1}{}>}
(перенаправлено с «P000401») может быть нулевой
в области, где <<<P000402> обозначение типа alias Declaration <<P000403>.
Мы говорим, что <<P000404>
<<<<P001480><<<P000405> Типовой псевдоним! Используется как класс
когда <<P000406> происходит одним из следующих способов:

<<P000502>
<<P001481>
<<P000407> происходит в случае создания выражения формы
<<P000725> или форму <<<P000726>,
где <<P000412> или \NEW{} или \CONST{}
(<<P000600>).
<<<P001484>{%>
Обратите внимание, что, например, <<<P000727> может быть творением,
Но затем он рассматривается как \code{\NEW{} $T$(42)} или <<P001487>>>{<<P001488>>>{} $T$(42)},
Что означает, что он включен здесь
(<<P000601>) %
?
<<<P001489>
<<<P000416> или <<P000728> является перенаправлением в перенаправлении заводского конструктора,
(<<P000602>).
<<P001490>
<<P000418> Используется в качестве суперкласса, смеси или суперинтерфейса.
в классовой декларации
(<<P000603>, <<P000604>, <<P000605>,
или как <<<P001491> Тип или суперинтерфейс в смешанной декларации
(<<P000606>).
% Его можно использовать как <<<P001492> Тип заявления о продлении:
% В этом положении он используется как тип, а не как класс.
<<P001493>
<<P000419> используется для вызова статического члена класса или смесителя,
Как описано выше.
<<P000526>

<<<P000877>%
Ошибка времени компиляции возникает, если <<<P000420> расширяется до переменной типа,
и <<<P000421> использует <<<P000422> как класс.
Ошибка времени компиляции возникает, если используется <<<P000423> как класс,
<<<P000425> не является регулярным.
<<<P001494> Например: }

<<P000002>

<<<P000878>%
Когда мы используем такие фразы, как Let <<<P000426> быть классом или предположить, что <<<P000427> Это смешение,
Следует понимать, что это включает случай, когда
<<P000428> <<<P001495>{typeName} обозначает псевдоним типа, или
<<P000429> Параметризованный тип формы
\syntax{<typeName> <typeArguments>>
где имя типа обозначает псевдоним типа,
и расширение транзитивного псевдонима <<<P000430> Класс обозначает соответственно смесь.


<<<P001497> Подтипы
<<P000899>

<<P000879>%
Этот раздел определяет, когда тип является <<<P000644> другого типа.
Ядром этого раздела является набор правил, определенных в
Рисунок ~<<<P000607>,
Но сначала нужно будет ввести несколько понятий.
Чтобы понять, что означают эти правила.

<<<P001498>{%>
Читатель, прочитавший множество научных работ об объектно-ориентированных системах
Смысл данного обозначения может быть очевиден,
Но нам все еще нужно прояснить некоторые детали о том, как справиться с этим.
синтаксически разные обозначения одного и того же типа,
Как выбрать правильную начальную среду (P000431).
%
Для читателя, который не знаком с обозначением, используемым в этом разделе,
Объяснений, приведенных здесь, должно быть достаточно, чтобы понять, что это значит.
Ссылаясь на объяснения естественного языка, приведенные в конце
раздел для получения интуиции о значении.%
?

<<<P000880>%
Этот раздел посвящен отношениям подтипов между типами.во время статического анализа
а также отношения подтипа, запрашиваемые в динамических проверках,
Типовые испытания
(<<P000608>),
и типовые отливки
(<<P000609>).

<<<P001499>{%>
Вариант правил, описанных здесь, показан в приложении.
(<<P000610>),
Демонстрация того, что субтипирование Dart может быть решено эффективно.
?

<<P000881>%
% % TODO(Eernst): Введение этих специализированных типов пересечений
% в подходящем месте, где указан тип продвижения.
Типы формы
\IndexCustom{$X \& S$>>>{тип!of} Форма <<P000433>%
\IndexExtraEntry{\&@$X \& S$>
возникают во время статического анализа из-за продвижения по типу
(P000611).
Они никогда не происходят во время казни.
Они никогда не являются аргументом другого типа.
ни тип возврата, ни тип формального параметра,
И это всегда так, что <<P000435> является подтипом границы <<<P000436>.
<<<P001502>{%>
Оригинальное название: P000437 является то, что он представляет
тип локальной переменной <<<P000438>
тип которого объявлен переменной типа <<<P000439>,
который, как известно, имеет тип <<<P000440> Из-за повышения.
Точно так же <<P000441> может рассматриваться как тип пересечения,
что является подтипом <<<P000442> А также подтип <<<P000443>.
Типы пересечений <<<P001503>{не} поддерживаются в целом.
Только в этом случае.%
?
Любая другая форма может возникнуть во время статического анализа.
Как и во время исполнения,
Взаимоотношения между подтипами всегда определяются одинаково.

% Номер правил подтипа
\newcommand{<<P001505>>>}{1}
\newcommand{<<P001507>>>}{2}
\newcommand{<<P001509>>>}{3}}
\newcommand{<<P001511>>>}{4}
\newcommand{<<P001513>>>}{5}
\newcommand{<<P001515>>>}{6}
\newcommand{<<P001517>>>}{7}
\newcommand{<<P001519>>>}{8]
\newcommand{<<P001521>>>}{9]
\newcommand{<<P001523>>>}{10]
\newcommand{<<P001525>>>}{11}
\newcommand{<<P001527>>>}{12]
\newcommand{<<P001529>>>}{13}
\newcommand{<<P001531>>>}{14}
\newcommand{<<P001533>>>}{15}
\newcommand{<<P001535>>>}{16}
\newcommand{<<P001537>>>}{17}
\newcommand{<<P001539>>>}{18}

<<<P000503>[p]
\def\VSP{\vspace{4 мм)
\def\ExtraVSP{\vspace{2 мм)
\def\Axiom#1#2#3#4{\centerline{\inference[#1]{}{<<P001550>>>{#3}{#4}}}<<<< P001551>>>
\def\Rule#1#2#3#4#5#6{\centerline{<<P001555>>>[#1]{<<P001556>>>{#3}{#4}}{<<P001557>>>{#5}{#6}}}<<<< P001558>>>
<<<P001559><<<P001560>#1#2#3#4#5#6#7#8
\centerline{<<P001562>>>[#1]{<<P001563>>>{#3}{#4} & \SubtypeStd{#5}{#6}}{<<P001565>>>{#7}{#8}}}\VSP
\def\RuleRaw#1#2#3#4#5
\centerline{<<P001570>>>[#1]{#3}{<<P001571>>>{#4}{#5}}}\VSP]
\def\RuleRawRaw#1#2#3#4{\centerline{\inference[#1]{#3}{#4}}\VSP]
%
<<P000504>>>[c]{0.49<<P001578>>>}
<<<P001579>{<<<P001580>> Рефлексивность.
\Axiom{\SrnBottom>>> Левый дно {<<<P001583>> T
<<P000527>
<<P000505>>>[c]{0.49<<P001584>>>}
<<<P001585>{<<<P001586>> Правый верх. <<P001587> <<<P001740>\code{Object} <<P001588>, \VOID\}}{S}{T}
\RuleRaw{\SrnNull>>> Левый нуль. <<<p001592>> \bot}{<<P000730>>>}{T}
<<P000528>

<<P001594>
<<<P001595>{<<<P001596>> Тип Alias Left {%}
\code{<<P001598>>>{) <<<P000444><\TypeParametersNoBounds{X}{}> = U} &
<<<<P001600>< S_1/X_1,\ldots,S_s/X_s]U}{T}}{\code{<<P000445>>><\List{S}{1}{}}}{T}
\RuleRaw{\SrnRightTypeAlias>>> Тип Alias Right {%}
\code{<<P001607>>>{) <<<P000446><\TypeParametersNoBounds{X}{}> = U} &
\SubtypeStd{S}{[T_1/X_1,<<P001610>>>,T_s/X_s]U}}{S}{<<<<<<<<<<<<<<<<<<<<<<<<<<<<<<<<<<<<<<<<<<<<<<<<<<<<<<<<<<<<<<<<<<<<<<<<<<<<<<<<<<<<<<<<<<<<<<<<<<<<<<<<<<<<<<<<<<<<<<<<<<<<<<<<<<<<<<<<<<<<<<<<<<<<<<<<<<<<<<<<<<<<<<<<<<<<<<<<<<<<<<<<<<<<<<<<<<<<<<<<<<<<<<<<<<<<< P001611>>>{$F$<\List T}{1}{}>}}

<<P000506>>>[c]{0.49<<P001613>>>}
<\RuleTwo{\SrnLeftFutureOr>>> Оставленное будущее {S}{T}{%}
\code{Future<<<<P000448>>>>}}{T}{<<P000732>>>}{T}\RuleTwo{\SrnRightPromotedVariable>>> Право Продвигаемая Переменная {S}{X}{S}{T}{%
S {\displaystyle X} <<P001742> T
<\Rule{\SrnRightFutureOrB>>> Правильное будущее Или B}{S}{T}{S}{\code{FutureOr<<<<P000450>>>>}
<\Rule{\SrnLeftVariableBound>>> Левая переменная граница {\Delta(X)}{T}{X}{T}
<<P000529>
<<P000507>>>[c]{0.49<<P001623>>>}
\Axiom{\SrnTypeVariableReflexivityA>>> Переменная A}{X <<<P001743> S}{X}
\Rule{\SrnRightFutureOrA>>> Правильное будущее Или A}{S}{<<P000734>>>}{%
S {<<<P000735>>
<\Rule{\SrnLeftPromotedVariable>>> Левая продвигаемая переменная B}{S}{T}{X <<<P001744> S {T }
\RuleRaw{\SrnRightFunction>>> Правильная функция {T<<<P001632> Тип функции {%}
T}{<<<P001633>>
<<P000530>
%
<<<P001634>
<<<<P001635>{<<<P001636>> Позиционные функции {%}
\Delta = \Delta\uplus\{X_ i\mapsto{}B_i<<<P001746><<<P001747>1 <<<P001641> i \leq s\} и
\Subtype{<<P001644>>>'}{S_0}{T_0}<<P001749>
n_1 <<<P001645> n_2
n_1 + k_1 <<<<P001646> n_2 + k_2
<<<P001647> j <<<P001648> 1.. n_2 + k_2\!:\;\Subtype{\Delta'} T_j} {S_j}}
<<<P000003>>
<<<P001651><<<P001652>
\RuleRawRaw{\SrnNamedFunctionType}{Названные типы функций}{%
\Delta = \Delta\uplus\{X_ i\mapsto{}B_i\, |\,1 <<<P001659> i \leq s\} и
\Subtype{<<P001662>>>}{S_0}{T_0} &
<<<P001663> j <<<P001664> 1.. n\!:\;\Subtype{\Delta'}{T_j}{S_j}<<P001758>>>
<<<P001759><<<P001760> \List{y}{n+1}{n+k_2}<<<<<<<<<<<<<<<<<<<<<<<<<<<<<<<<<<<<<<<<<<<<<<<<<<<<<<<<<<<<<<<<<<<<<<<<<<<<<<<<<<<<<<<<<<<<<<<<<<<<<<<<<<<<<<<<<<<<<<<<<<<<<<<<<<<<<<<<<<<<<<<<<<<<<<<<<<<<<<<<<<<<<<<<<<<<<<<<<<<<<<<<<<<<<<<<<<<<<<<<<<<<<<<<<<<<<<<<<<<<<<<<<<<<<<<<< P001761>>><<<P001762> <<<P001668> \{\,\List{x}{n+1}{n+k_1}\,<<<P001766> <<P001767>
<<<P001670 >> p <<<P001671> 1.. k_2, q \in 1.. k_1:\quad
y_{n+p} = x_{n+q}\quad\Rightarrow\quad\Subtype{\Delta>{T_{n+p}}{S_{n+q}}}{%
<<<P000004>>
%
<<<P001679>
\RuleRaw{\SrnCovariance>>> Ковариантность
\code{<<P001683>>>{} $C$\TypeParametersNoBounds{X}{}><<P001768>>><<P001685>>>\,\{\} и
<<<P001686> j <<<P001687> 1..s\!:\;\SubtypeStd{S_j}{T_j}}{%
<<<P001689><<<P000454><<<<P001690> S}{1}{}>}}{<<< P001691>>>{$C$<<<<P001692> T}{1}{}>}}
<<<P001693>
<\RuleRaw{\SrnSuperinterface>>> Суперинтерфейс (%)
\code{<<P001697>>>{) $C$\TypeParametersNoBounds{X}{}>\,\ldots\,\{\}\\
\Superinterface{<<P001701>>>{<<P000457>>><\List{T}{1}{m}> и
<<<<P001703>< S_1/X_1,\ldots,S_s/X_s]\code{<<P000458>>><\List{T}{1}{m}>}}
<<<P001707><<<P000459><<<<P001708> S}{1}{}>}\;}{<<P001780>>> T
%
<<<P001709> Правила подтипа.
<<P000613>
<<P000531>


<\subsubsection{{{{{{{{{{{{{{{{{{{{{{{{{{{{{{{{{{{{{{{{{{{{{{{{{{{{{{{{{{{{{{{{{{{{{{{{{{{{{{{{{{{{{{{{{{{{{{{{{{{{{{{{{{{{{{{{{{{{{{{{{}}{{{{{{{{{{{{{{{{{{{{{{{{{{{{{{{{{{{{{{{{{{{{{{{{{{{{{{{{{{{{{{{{{{{{{{{{{{{{{{{{{{{{{{{{{{{{{{{{{{{{{{{{{{{{{{{{{{{{{{{{{{{{{ Мета-варианты.
<<P000900>

<<P000882>%
<<<P000645> Символ, обозначающий синтаксическую конструкцию
удовлетворяет некоторым статическим семантическим требованиям.

<<<P001711>{%>
Например, <<<P000460> является мета-вариабельным, обозначающим
идентификатор <<<P000736>,
Но только если <<P000737> обозначает переменную типа, объявленную в замкнутом пространстве.
В приведенных ниже определениях мы определяем это, говоря, что
<<P000461> Диапазоны по переменным типа.
Точно так же <<P000462> является мета-вариабельным, обозначающим
a \synt{typeName}, например, <<P000738>,
Но только если <<P000739> обозначает класс в заданном объеме.
Мы определяем это как <<<P000463> Диапазон классов.%
?

<<P000883>%
В этом разделе мы используем следующие мета-варианты:

<<P000508>
<<<P001713> <<P000464> Диапазоны по переменным типа.
<<<P001714> <<P000465> Диапазон классов,
<<<P001715> <<P000466> Диапазоны по типам псевдонимов.<<<P001716> <<<P000467> и <<P000468> диапазон по типам, возможно с индексом, таким как <<P000469> или <<P000470>.
<<<P001717> <<P000471> диапазоны по типам, опять же, возможно, с индексом;
Используется только как переменная типа.
<<P000532>


<<<P001718> Правила подтипа
<<P000901>

<<P000884>%
В этом разделе мы определяем несколько правил о субтитрах.
Всякий раз, когда правило содержит одну или несколько метавариабель,
Это правило может быть использовано
\IndexCustom{instantiating}{instantiation!subtype rule}
то есть путем последовательной замены
Каждое появление данной метавариабельности
конкретный синтаксис, обозначающий один и тот же тип.

<<<P001720>{%>
Это означает, что два или более случая
заданная мета-вариабельность в правиле
обозначает идентичные части синтаксиса,
и введение правила продолжается как
Простая операция поиска и замены.
Например,
Правило ~\SrnReflexivity На рисунке ~<<<P000612>
Можно использовать для заключения
\Subtype{<<P001723>>>}{<<P000740>>>}{<<P000741>>>}
где <<P000472> обозначает пустую среду
(Любой среды будет достаточно, потому что переменные типа не встречаются).

Однако вышеприведенная формулировка "обозначение одного и того же типа" охватывает
Дополнительные ситуации также:
Например, мы можем использовать правило ~\SrnReflexivity
Чтобы показать, что <<P000742> является подтипом
<<P000743>, когда <<P000744> Класс, объявленный в
библиотека <<P000473> которые импортируются библиотеками <<<P000474> и <<<P000475> и
используется в декларациях,
<<<P000476> и <<<P000477> Импортируются с префиксами
<<P000745> <<<P000746> по текущей библиотеке.
Важно отметить, что все случаи одной и той же мета-вариабельности
в данном правиле инстанциация обозначает один и тот же тип,
Даже в том случае, если этот тип не обозначен
Один и тот же синтаксис в обоих случаях.

И наоборот, мы можем использовать одно и то же правило, чтобы сделать вывод.
<<<P000747> является подтипом <<<P000748>
в случае, когда первый обозначает класс, объявленный в библиотеке <<<P000478>
и последний обозначает класс, объявленный в <<<P000479>, с <<<P000480>.
Эта ситуация может возникнуть без ошибок времени компиляции, например,
Если <<<P000481> и <<<P000482> Косвенный ввоз в нынешнюю библиотеку
и два «смысла» <<<P000749> используются
аннотации типа на переменные или формальные параметры функций
задекларировано в промежуточных библиотеках импортом <<<P000483> соответственно <<P000484>.
Неспособность доказать
\Subtype{<<P001727>>>}{<<P000750>>>}{<<P000751>>>] " "
Это произойдет, например, в ситуации, когда мы проверим,
такая переменная может быть передана в качестве фактического аргумента такой функции;
Потому что два случая <<P000752> не обозначают один и тот же тип.%
?

